\section{Worked finite examples}
\label{sec:examples}

The accompanying Python program is an exact reference implementation using
integers and \code{fractions.Fraction}.  It is not a proof in \(\ipra\), but
its examples mirror the finite recurrences to be embedded.

\subsection{A bounded trace}

Fix an input \(u\) and define
\[
 s_0=0,\qquad s_{i+1}=s_i+u+i.
\]
The invariant is
\[
 s_i=iu+\frac{i(i-1)}2.                         \tag{E1}\label{eq:trace-demo}
\]
The base case is immediate.  At successor,
\[
 s_{i+1}=iu+\frac{i(i-1)}2+u+i
 =(i+1)u+\frac{i(i+1)}2.
\]
Because \(i(i-1)\) is even, one may either use coded exact division by two or
avoid division by storing \(2s_i=2iu+i(i-1)\).  The latter is the cleaner
quantifier-free invariant.

For \(u=3\) and bound \(8\), the certified states are
\[
 0,3,7,12,18,25,33,42,52.
\]
The checker validates every local step and every instance of
\eqref{eq:trace-demo}.  The final value is \(52\).

\subsection{An overlapping decision}

Take
\[
 a=\frac9{16},\qquad e=\frac1{32},\qquad
 q=\frac12,\qquad d=\frac18.
\]
Here \(4e=d\), and
\(a=q+d/2\).  The upper branch of \cref{thm:overlap-branch} returns
\[
 x<q+d=\frac58.
\]
Indeed \(x\le a+e=19/32<20/32=5/8\).  Had the approximation been larger than
\(9/16\), the lower branch would certify \(x>1/2\).  The interval
\((1/2,5/8)\) is deliberately shared by both branches.

This tiny example is the complete control-flow core of much larger chart and
pivot arguments.

\subsection{Finite-grid approximation to \texorpdfstring{\(\sqrt2\)}{sqrt(2)}}

Let \(P(x)=x^2-2\) on \([1,2]\).  On this interval,
\cref{lem:poly-lip} gives Lipschitz constant \(4\).  Scan the grid with
denominator \(16\).  The least point minimizing \(\abs{P}\) is
\[
 x_* = \frac{23}{16},
 \qquad
 \abs{P(x_*)}=\frac{17}{256}.
\]
A nearest grid point to \(\sqrt2\) is within \(1/32\), so the generic net
bound gives residual at most \(4/32=1/8\); the exact residual above is smaller.
The scan needs no square-root function.  It uses only integer squares and
rational comparison.

To obtain an exact regular name, at precision \(k\) compute
\[
 n_k=\floor{2^{k}\sqrt2}
     =\mathsf{isqrt}(2\cdot2^{2k}),
\]
checked by
\[
 n_k^2\le2\cdot2^{2k}<(n_k+1)^2.
\]
The interval
\([n_k/2^k,(n_k+1)/2^k]\) has width \(2^{-k}\), contains \(\sqrt2\), and is
nested up to a harmless index shift.  Integer square root is a bounded search
with a quantifier-free invariant.  The archive's regression test checks the
resulting midpoint regularity through precision twelve; the general proof is
the nested-ball theorem.

\subsection{Polynomial identity certificates}

Consider
\[
 (x+y)^2=x^2+2xy+y^2.
\]
A certificate producer expands the left side, sorts monomials
\((2,0),(1,1),(0,2)\), and returns coefficients \((1,2,1)\).  The right side
already has that normal form.  The checker verifies each list transformation.
Soundness is proved once by induction on the certificate list; the theorem for
this identity is then one accepted finite certificate plus universal
introduction of \(x,y\).

For a rational identity, the checker also records a positive common
denominator.  The final statement is the corresponding integer polynomial
identity after multiplication.  This is the recommended path for matrix,
trigonometric, and algebraic identities in the packets.

\subsection{Height bound and list}

Let \(P(a,x)\) be \(x^2\le a\), with height \(h(x)=x\).  For \(a=20\), the
height certificate \(H=4\) produces the complete list
\([0,1,2,3,4]\).  Conversely the maximum of the list is \(4\).  The example is
trivial, but its proof structure is identical to bounded rational-point search:
a theorem-specific height implication feeds a generic finite filter.

\subsection{Finite probability}

Let a law have rational values \(x_j=a_j/b_j\) and natural weights \(w_j\),
with \(W=\sum_jw_j>0\).  Its expectation is the exact rational
\[
 \mathbb E X=\frac{1}{W}\sum_jw_jx_j.
\]
For two laws, enumerate pairs by mixed radix.  The identity
\[
 \sum_{i,j}w_i v_j f(x_i)g(y_j)
 =\left(\sum_iw_i f(x_i)\right)
  \left(\sum_jv_j g(y_j)\right)
\]
is a finite polynomial certificate.  It proves independence algebra, mean
additivity, and variance addition for centered variables.  There is no measure
space and no appeal to Fubini; all sums are finite and their indices have
explicit bounds.

\subsection{Reproduction}

From the archive root, run
\begin{verbatim}
cd code
python3 -m unittest -v test_certificate_kernel.py
python3 certificate_kernel.py ../examples/framework_certificate.json
\end{verbatim}
The delivered run passes nine tests.  The JSON certificate records the trace,
branch, finite-grid, and regularity examples.  These tests provide independent
finite arithmetic checks; they do not replace the planned \code{Prf PRAAx}
formalization.
